\documentclass[aspectratio=169, ignorenonframetext]{beamer}
\usepackage{pgf}

\mode<article>{
  \usepackage{fullpage}

}

\mode<presentation>
{
%   \usetheme{Madrid}
%   \useinnertheme{circles}
  \setbeamertemplate{navigation symbols}{}
%   \usetheme[hideothersubsections, width=2.4cm]{Hannover}
%   \usetheme{Antibes}
%   \usetheme{Montpellier}
%   \usetheme{Singapore}
  \usetheme{Darmstadt}
%   \usecolortheme{seahorse}
%   \usecolortheme{orchid}
  \setbeamertemplate{footline}
  {%
%     \begin{beamercolorbox}{section in head/foot}
%       \usebeamercolor{bg}
       \hskip 1em \footnotesize \insertframenumber{} / \inserttotalframenumber%\hskip 41em \includegraphics[height = .5cm]{../../../../bilder/cc_by-nc_eu.png}
% % cc_by-nc_eu.png: 403x141 px, 72dpi, 14.22x4.97 cm, bb=0 0 403 141
%
% % bwslogo_3.png: 476x392 px, 300dpi, 4.03x3.32 cm, bb=0 0 114 94
%
%       %\hskip 5em
%       %       \input{../bilder/cc_by.png}
%       %\includegraphics[height=.5cm]{../../../../bilder/bwslogo_3.png}
%     \end{beamercolorbox}%
  }
  \usepackage{beamerfoils}
}
\usepackage[german]{babel}
\usepackage[utf8]{inputenc}
\usepackage{times}
\usepackage[T1]{fontenc}
\usepackage{pgf}
\usepackage{pgfpages}
\usepackage{eurosym}
\usepackage{graphicx}
\usepackage{amsmath}
\usepackage[siunitx,european]{circuitikz}
\usepackage{ulem}
\usepackage{listings}
\usepackage{svg}
%
\lstset{numbers=left, numberstyle=\tiny, stepnumber=2, numbersep=5pt, language = C++, alsolanguage=XML}
% \MyLogo{\includegraphics[height=1cm]{../../../../bilder/bwslogo_3.png}}
% % \includegraphics{../../bilder/bwslogo_3.png}
% % bwslogo_3.png: 476x392 px, 300dpi, 4.03x3.32 cm, bb=
%
\only<article>{
  \usepackage[colorlinks=true,linkcolor=blue,filecolor=magenta,urlcolor=cyan]{hyperref}
}

\only<presentation>{
  \usepackage{hyperref}
}


\title{Arbeitsunterlagen zu FOS Elektrotechnik, technische Informatik, Mechatronik Themenfeld 12.1}
\subtitle{Gleichstromnetzanalyse - Teil 1}
\date{V 0.5.1 \\ Stand: \today}%\\

\institute[BWS Hofheim]{Brühlwiesenschule, Hofheim}
\author{Thomas Maul}

\titlegraphic{Für eigene Teile gilt: \includegraphics[height=1cm]{cc_by-nc_eu.png}}

\begin{document}
  \only<article>{
    \maketitle
    \tableofcontents
    \clearpage
  }
  \begin{frame}<beamer>
    \titlepage
    % \hyperlink{Teil_2}{\beamerbutton{Go part 2}}
  \end{frame}
  % \AtBeginSection[] % Do nothing for \section*
  % {
  %   \begin{frame}<beamer>
  %     \frametitle{Inhalt}
  %     \tableofcontents[currentsection]
  %   \end{frame}
  % }
% \begin{frame}{Was bisher geschah}
%     \begin{itemize}
%         \item Atome und Elektronen
%         \item Ladungen im Sinne der Elektrotechnik
%         \item Gleichstrom und Wechselstrom
%         \item Widerstand, Strom, Spannung, Leistung
%     \end{itemize}
% \end{frame}


% \part{Themenfeld 12.1 - Gleichstromnetzanalyse}
\begin{frame}
  \partpage
  %   \tableofcontents[hidesubsections]
\end{frame}
\begin{frame}<beamer>
  \frametitle{Inhalt}
  \begin{columns}
    \column{.5\textwidth}
    \tableofcontents[sections={1-4},hidesubsections]%currentsection]
    ZAP = In der Themenliste der Abschlussprüfung genannt.
    \column{.5\textwidth}
    \tableofcontents[sections={5-},hidesubsections]%currentsection]
  \end{columns}
\end{frame}
\section{Ersatzquellen}

\begin{frame}{Umwandlung Stromquelle $\rightarrow$ Spannungsquelle I}
 \begin{columns}
   \begin{column}{.5\textwidth}
      Ideale Stromquelle
     \begin{itemize}
     \item liefert definierten, konstanten Strom
     \item Spannung ist unerheblich
     \item Innenwiderstand ($R_i$) ist parallel geschaltet.
    \end{itemize}
   \end{column}
   \begin{column}{.5\textwidth}
    \begin{figure}[htb]
      \begin{circuitikz}
        \draw (0,0) to[isource, l=$I_{q1}$] (0,2) -- (1,2);
        \draw (1,2) -- (2,2);
        \draw (1,2) node[circ]{} -- (2,2) node[ocirc]{} node[above]{A} ;
        \draw (1,2) to[R=$R_i$] (1,0) node[circ]{} -- (0,0);
        \draw (2,2) to[R=$R_1$] (4,2) -- (4,0)  -- (0,0);
        \draw (2,0) node[ocirc]{} node[below]{B};
%         \draw (2,2) -- (3,2) node[ocirc]{} node[above]{C} -- (3.5,2) to[R=$R_3$]
%         (6,0) -- (3,0) node[ocirc]{} node[below]{D} -- (2,0);
%         \draw (3.5,2) -- (4,2) to[R=$R_4$] (6,2);
%         \draw (3.5,2) node[circ]{} (6,0) node[circ]{};
%         \draw (6,2) to[R=$R_5$] (6,0);
      \end{circuitikz}
      \caption{Stromquelle}
    \end{figure}
    \label{fig:Stromquelle1}
   \end{column}
 \end{columns}
\end{frame}

\begin{frame}{Umwandlung Stromquelle $\rightarrow$ Spannungsquelle II}
 \begin{columns}
   \begin{column}{.5\textwidth}
      Stromquelle in Spannungsquelle
     \begin{itemize}
     \item Wert von $R_i$ bleibt,
     \item $R_i$ geht in Reihe zu $U_q$
     \item Spannung: $U_l = I_k \cdot R_i$
    \end{itemize}
   \end{column}
   \begin{column}{.5\textwidth}
    \begin{figure}[htb]
      \begin{circuitikz}
        \draw (0,0) to[vsource, l=$I_{q1}$] (0,2) ;
        \draw (3,2) node[ocirc]{} node[above]{A} ;
        \draw (0,2) to[R=$R_i$] (2,2) -- (3,2);
        \draw (3,2) to[R=$R_1$] (5,2) -- (5,0)  -- (0,0);
        \draw (3,0) node[ocirc]{} node[below]{B};
%         \draw (2,2) -- (3,2) node[ocirc]{} node[above]{C} -- (3.5,2) to[R=$R_3$]
%         (6,0) -- (3,0) node[ocirc]{} node[below]{D} -- (2,0);
%         \draw (3.5,2) -- (4,2) to[R=$R_4$] (6,2);
%         \draw (3.5,2) node[circ]{} (6,0) node[circ]{};
%         \draw (6,2) to[R=$R_5$] (6,0);
      \end{circuitikz}
      \caption{Spannungsquelle}
    \end{figure}
    \label{fig:Spannungsquelle1}
   \end{column}
 \end{columns}
\end{frame}

\begin{frame}{Umwandlung Spannungsquelle $\rightarrow$ Stromquelle}
 \begin{columns}
   \begin{column}{.5\textwidth}
      Spannungsquelle in Stromquelle
     \begin{itemize}
     \item Wert von $R_i$ bleibt
     \item $R_i$ ist jetzt parallel zu $I_q$
     \item Strom: $I_k = \frac{U_l}{R_i}$
    \end{itemize}
   \end{column}
   \begin{column}{.5\textwidth}
    \begin{figure}[htb]
      \begin{circuitikz}
        \draw (0,0) to[isource, l=$I_{q1}$] (0,2) -- (1,2);
        \draw (1,2) -- (2,2);
        \draw (1,2) node[circ]{} -- (2,2) node[ocirc]{} node[above]{A} ;
        \draw (1,2) to[R=$R_i$] (1,0) node[circ]{} -- (0,0);
        \draw (2,2) to[R=$R_1$] (4,2) -- (4,0)  -- (0,0);
        \draw (2,0) node[ocirc]{} node[below]{B};
%         \draw (2,2) -- (3,2) node[ocirc]{} node[above]{C} -- (3.5,2) to[R=$R_3$]
%         (6,0) -- (3,0) node[ocirc]{} node[below]{D} -- (2,0);
%         \draw (3.5,2) -- (4,2) to[R=$R_4$] (6,2);
%         \draw (3.5,2) node[circ]{} (6,0) node[circ]{};
%         \draw (6,2) to[R=$R_5$] (6,0);
      \end{circuitikz}
      \caption{Stromquelle}
    \end{figure}
    \label{fig:Stromquelle2}
   \end{column}
 \end{columns}
\end{frame}

\begin{frame}{Stromquelle mit "`zusätzlichem"' Widerstand}
 \begin{columns}
   \begin{column}{.3\textwidth}
      Ideale Stromquelle
     \begin{itemize}
     \item $R_i$ ist parallel zu $I_q$
     \only<2>{\item R2 ist irrelevant}
    \end{itemize}
   \end{column}
   \begin{column}{.7\textwidth}
    \begin{figure}[htb]
      \begin{circuitikz}
        \draw (0,0) to[isource, l=$I_{q1}$] (0,2) ;
        \draw (0,2) to[R=$R_2$] (2,2);
        \draw (2,2) node[circ]{} -- (3,2) node[ocirc]{} node[above]{A} ;
        \draw (2,2) to[R=$R_3$] (2,0) node[circ]{} -- (0,0);
        \draw (3,2) to[R=$R_1$] (5,2) -- (5,0)  -- (0,0);
        \draw (3,0) node[ocirc]{} node[below]{B};
%         \draw (2,2) -- (3,2) node[ocirc]{} node[above]{C} -- (3.5,2) to[R=$R_3$]
%         (6,0) -- (3,0) node[ocirc]{} node[below]{D} -- (2,0);
%         \draw (3.5,2) -- (4,2) to[R=$R_4$] (6,2);
%         \draw (3.5,2) node[circ]{} (6,0) node[circ]{};
%         \draw (6,2) to[R=$R_5$] (6,0);
      \end{circuitikz}
      \caption{Stromquelle}
    \end{figure}
    \label{fig:StromquelleMitZusR}
   \end{column}
 \end{columns}
\end{frame}

\begin{frame}{Spannungsquelle mit "`zusätzlichem"' Widerstand}
 \begin{columns}
   \begin{column}{.3\textwidth}
      Ideale Spannungsquelle
     \begin{itemize}
     \item $R_i$ ist in Reihe zu $U_q$
     \only<2>{\item R3 ist irrelevant}
    \end{itemize}
   \end{column}
   \begin{column}{.7\textwidth}
    \begin{figure}[htb]
      \begin{circuitikz}
        \draw (0,0) to[vsource, l=$U_{q1}$] (0,2) -- (1,2);
        \draw (1,2) to[R=$R_2$] (3,2);
        \draw (1,2) node[circ]{};
        \draw (3,2) node[ocirc]{} node[above]{A} ;
        \draw (1,2) to[R=$R_3$] (1,0) node[circ]{} -- (0,0);
        \draw (3,2) to[R=$R_1$] (5,2) -- (5,0)  -- (0,0);
        \draw (3,0) node[ocirc]{} node[below]{B};
%         \draw (2,2) -- (3,2) node[ocirc]{} node[above]{C} -- (3.5,2) to[R=$R_3$]
%         (6,0) -- (3,0) node[ocirc]{} node[below]{D} -- (2,0);
%         \draw (3.5,2) -- (4,2) to[R=$R_4$] (6,2);
%         \draw (3.5,2) node[circ]{} (6,0) node[circ]{};
%         \draw (6,2) to[R=$R_5$] (6,0);
      \end{circuitikz}
      \caption{Spannungsquelle}
    \end{figure}
    \label{fig:SpannungsquelleMitZusR}
   \end{column}
 \end{columns}
\end{frame}
\end{document}
